\documentclass[10pt,a4paper]{amsart}
\usepackage[margin=19mm]{geometry}
\usepackage{amsmath,amssymb,mathtools}
\usepackage[scr=rsfs,cal=euler]{mathalfa}
\usepackage{xfrac}
\usepackage[unicode,hidelinks,bookmarks=false]{hyperref}
\newcommand{\ie}{i.e.\ }
\newcommand{\cf}{cf.\ }
\newcommand{\resp}{respectively\ }
\newcommand{\Subsection}{\subsection}
\providecommand{\nobreakdash}{\nobreak\hbox{-}}
\setlength{\emergencystretch}{2em}
\multlinegap0pt
\usepackage{amssymb}
\usepackage{tikz}
\newtheorem{theorem}{Theorem}
\newcommand{\GGnn}[1]{{#1}}
\newcommand{\wdom}[1]{\rightsquigarrow}
\title{Forward Arc Maximization for Hamilton Oriented Cycles and Paths in Generalizations of Tournaments}
\author{Q. Guo, G. Gutin, Y. Lan, Q. Shao, A. Yeo, Y. Zhou}
\date{Source: arXiv:2602.10713v1 (CC BY 4.0)}
\begin{document}
\maketitle

\noindent\emph{Source and licence.} Forward Arc Maximization for Hamilton Oriented Cycles and Paths in Generalizations of Tournaments. Version arXiv:2602.10713v1 (CC BY 4.0), distributed by arXiv under the Creative Commons Attribution 4.0 International licence, http://creativecommons.org/licenses/by/4.0/, as recorded in the arXiv licence metadata for this exact version and shown on the abstract page https://arxiv.org/abs/2602.10713v1. Full licence notice: \texttt{licenses/arxiv-2602.10713-CC-BY-4.0.txt}.\par\smallskip

\noindent\textit{Editorial scope.} Counted unit: the article's theorem main\_irreducible\_simple (source0.tex lines 411--416). The article's complete original proof of it is delivered with it (source0.tex lines 428--464, one \texttt{proof} environment of 37 lines, which the article places away from the statement). No mathematics is written, repaired or re-derived: the statement and the proof are the article's own lines, in the article's own notation, and no retained line is substituted or re-flowed.  No further source-local context is needed: every cross-reference of every retained line resolves inside the two delivered spans.  Nothing was omitted from any retained span, so the package carries no omission record.  No retained line contains a citation, so the wrapper carries no bibliography and no bibliography entry.  No retained line contains a figure environment, an image-inclusion command or drawing code, so no image is delivered and the package declares no asset dependency.  Editorial formatting only, in addition: an AMS \texttt{amsart} preamble that restates the article's own declarations and macros used by the retained lines, namely the environments \texttt{theorem} on separate counters, together with the standard packages those lines need.  The retained environments print in this document as Theorem 1 in order of appearance, whereas the article prints the same items with the numbering of its own full document.  The complete native source of the article as distributed in the arXiv e-print for this exact version is bundled unchanged as \texttt{sources/194421/source0.tex}.
\begin{theorem}\label{main_irreducible_simple}
Let $D$ be a semicomplete multipartite digraph. \GGnn{In polynomial time, we can decide whether $D$ has a cycle factor and if it has one, then, in polynomial time, we can find either a Hamilton cycle  or a cycle factor with at least $t\ge 2$ cycles in which we can order its cycles $C_1,C_2,\ldots, C_t$ such that 
 %$F$, and if $F$ is not a Hamilton cycle, then its cycles can be ordered 
%one can find a cycle factor $F$ with $t$ cycles of $D$, and if $t\ge 2$, an ordering $C_1,C_2,\ldots, C_t$ of the cycles of $F$ which satisfies the following property:
$C_i \wdom{} C_j$ for all $1 \leq i < j \leq t$.}
\end{theorem}

\begin{proof}

Let $D$, ${F}$ and $P$ be defined as in the statement of the theorem. Let $P=p_1 p_2 p_3 \ldots p_{\ell}$ and
let $D'$ be obtained from $D$ by adding the arc $p_{\ell} p_1$ if it does not exist in $D$ already (if $p_{\ell} p_1 \in A(D)$ then $D'=D$).
Note that $D'$ has a cycle factor. 
%and let ${C}$ denote a cycle factor in $D'$.
%If the arc $p_{\ell} p_1$ belongs to a cycle in ${C}$ then let $a=p_{\ell} p_1$ and otherwise let $a$ be an arbitrary arc on a cycle in ${C}$.

%If ${C}$ contains only one cycle then we obtain the desired Hamilton path by removing the arc $a$ from the cycle.
%So we may assume that ${ C}$ contains at least two cycles.
%We can therefore use 
By Theorem~\ref{main_irreducible_simple}, we can find in polynomial time a cycle factor $C$  of $D'$ with $t$ cycles and either $t=1$ or we can order its cycles, $C_1,C_2,\ldots,C_t$, of ${C}$ such that
$C_i \wdom{} C_j$ for all $1 \leq i < j \leq t$, where $t \geq 2$. If the arc $p_{\ell} p_1$ belongs to a cycle in ${C}$ then let $a=p_{\ell} p_1$ and otherwise let $a$ be an arbitrary arc on a cycle in ${C}$.
Assume that the arc $a$ belongs to $C_r$, where $1 \leq r \leq t$ and consider the $1$-path-cycle factor, ${F}'$, in $D$ obtained from
${C}$ by deleting the arc $a$. Let $P'=p_1' p_2' p_3' \ldots p_m'$ be the path in ${F}'$.

Note that $p_1'$ and $p_m'$ belong to different partite sets, and assume without loss of generality that $p_1' \in V_1$ and $p_m' \in V_2$, where
$V_1,V_2,\ldots,V_c$ are the partite sets of $D$.

First consider the case when $r<t$. We now transform $P'$ and $C_{r+1}$ into a path where the end-points belong to different partite sets as follows.
We first consider the case when there is no arc from $C_{r+1}$ to $p_m'$ in $D$. In this case let $y \in V(C_{r+1}) \cap V_1$ if such a vertex exists and otherwise let $y \in V(C_{r+1}) \setminus V_2$ be arbitrary. Note that $p_m' y \in A(D)$ and the path $P' C_{r+1}[y,y^-]$ is a path with vertex set $V(C_r) \cup V(C_{r+1})$ and with end-points in different partite sets.

Secondly we consider the case when there is an arc from $C_{r+1}$ to $p_m'$ in $D$, say $z p_m'$.
By Theorem~\ref{main_irreducible_simple} we note that $p_{m-1}'$ and $z^+$ both belong to the same partite set $V_i$ \GGnn{and $p'_{m-1}z,p'_mz^+\in A(D).$} 
Hence, the following are both paths on the vertex set $V(C_r) \cup V(C_{r+1})$:
\[
P_1 = p_1' p_2' \cdots p_{m-1}' z p'_m C_{r+1}[z^+,z^-] \mbox{ and }
P_2 = p_1' p_2' \cdots p_{m}' C_{r+1}[z^+,z]
\]
And as $z$ and $z^-$ belong to different partite sets either $P_1$ or $P_2$ has initial and terminal vertices from different partite sets. So we have in all cases
found a path with vertex set $V(C_r) \cup V(C_{r+1})$ with  initial and terminal vertices in different partite sets.

We can repeat this process in order to get a path with vertex set $V(C_r) \cup V(C_{r+1}) \cup V(C_{r+2})$ with  initial and terminal vertices in different partite sets.
And continuing this process further we obtain a path with vertex set $V(C_r) \cup V(C_{r+1}) \cup \cdots \cup V(C_t)$ with  initial and terminal vertices in different partite sets.

We can now analogously merge cycles $C_{r-1}$, $C_{r-2}$, $\ldots$, $C_1$ with the path ending up with a Hamilton path in $D$ where the initial and terminal vertices are in different partite sets.
\end{proof}

\end{document}
